\section{Conclusion}
\label{sec:conclusion}

Penelitian ini membuktikan efektivitas fusi representasi multi-domain ViT (Face $\oplus$ Emotion $\oplus$ Age: 2.304 dimensi) untuk klasifikasi enam kelas interseksional dari tiga kelompok ras makro (Asian, Black, White) pada dataset DemogPairs. Di bawah protokol Strict Subject-Disjoint Partitioning dengan 0\% kebocoran identitas (480 subjek latih dan 120 subjek uji held-out), pengklasifikasi LR meraih performa puncak dengan akurasi 88.66\%, presisi makro 88.75\%, recall makro 88.66\%, dan F1-score makro 88.68\%, disusul pengklasifikasi SVM dengan akurasi 88.61\%, presisi makro 88.71\%, recall makro 88.61\%, dan F1-score makro 88.63\%. Pengujian hipotesis McNemar berpasangan membuktikan signifikansi fusi tri-domain terhadap representasi domain tunggal Face ($p < 0.001$) serta mengonfirmasi paritas prediktif antara LR dan SVM ($p = 1.0000$, dengan selisih bersih satu sampel dari 2.160 citra). Evaluasi keadilan algoritmik antarsubkelompok menunjukkan stabilitas prediksi dengan rentang recall 84.72\% hingga 94.44\% serta kesenjangan EOD sebesar 9.72~pp pada kedua model teratas.

Meskipun menawarkan kontribusi empiris yang signifikan, penelitian ini memiliki dua batasan utama yang memerlukan perhatian. Pertama, cakupan demografis penyelidikan ini terbatas pada tiga kelompok ras makro dalam dataset DemogPairs (Asian, Black, White), sehingga generalisasi performa klasifikasi ke kelompok multirasial maupun populasi etnisitas lain tetap membutuhkan pengujian empiris tambahan. Kedua, proses ekstraksi representasi multi-domain masih dijalankan secara terpisah (offline) memanfaatkan tiga model backbone ViT-Base beku tanpa fine-tuning end-to-end terpadu, sehingga memerlukan tahapan pemrosesan fitur tersendiri sebelum klasifikasi hilir. Menanggapi batasan tersebut, arah riset mendatang diarahkan pada tiga prioritas strategis. Pertama, validasi lintas dataset (cross-dataset validation) perlu dilakukan pada korpus eksternal berskala besar in-the-wild, seperti FairFace atau UTKFace, untuk menguji ketangguhan model terhadap variasi pose ekstrem, oklusi, dan pencahayaan dinamis. Kedua, penelitian lanjutan perlu mengembangkan mekanisme fusi atensi silang adaptif antardomain (end-to-end cross-attention fusion) serta menerapkan teknik kuantisasi atau knowledge distillation. Ketiga, integrasi metode visualisasi Explainable Artificial Intelligence (XAI), seperti Attention Rollout, menjadi langkah krusial untuk menyediakan transparansi representasi fitur serta mendukung prinsip Responsible AI.

\section*{Data Availability Statement}
Dataset DemogPairs yang dianalisis dalam penelitian ini merupakan dataset yang tersedia secara publik untuk keperluan riset akademis. Citra wajah beserta anotasi label demografis dapat diakses secara bebas melalui publikasi resmi pengembang pada \url{https://ihupont.github.io/publications/2019-05-16-demogpairs} dan direktori penyimpanan Google Drive resmi pada \url{https://drive.google.com/file/d/1f_ez-ll6wxDXScrG4ceStZRuLT_K8Uy8/view?usp=sharing}.
